\documentclass[a4paper, 10pt, twocolumn]{article}

% --- Packages ---
\usepackage[utf8]{inputenc}
\usepackage{graphicx}
\usepackage{booktabs}   % For professional table formatting
\usepackage{tabularx}   % For adjustable width tables
\usepackage{titlesec}   % For formatting section headings
\usepackage{geometry}   % For setting margins
\usepackage{hyperref}   % For clickable links
\usepackage{lipsum}     % For generating Latin filler text

% --- Page Setup ---
\geometry{
    top=2cm,
    bottom=2cm,
    left=1.5cm,
    right=1.5cm,
    columnsep=0.8cm
}

% Format section headings to match the source document
\titleformat{\section}{\normalfont\large\bfseries\uppercase}{\thesection.}{1em}{}
\titleformat{\subsection}{\normalfont\normalsize\bfseries}{\thesubsection}{1em}{}

% --- Title and Authors ---
\title{\textbf{LOREM IPSUM DOLOR SIT AMET CONSECTETUR ADIPISCING ELIT}}
\author{
    First Author$^{1}$, Second Author$^{2}$ \\
    \small $^{1}$ First Department, University of Example \\
    \small $^{2}$ Second Department, Institute of Testing \\
    \small \{author.one\}@example.com, \{author.two\}@example.ac.uk
}
\date{} % Leave empty to hide the date

\begin{document}

\maketitle

% --- Abstract ---
\begin{abstract}
Lorem ipsum dolor sit amet, consectetur adipiscing elit. Sed do eiusmod tempor incididunt ut labore et dolore magna aliqua. Ut enim ad minim veniam, quis nostrud exercitation ullamco laboris nisi ut aliquip ex ea commodo consequat. Duis aute irure dolor in reprehenderit in voluptate velit esse cillum dolore eu fugiat nulla pariatur. Excepteur sint occaecat cupidatat non proident, sunt in culpa qui officia deserunt mollit anim id est laborum.
\end{abstract}

% --- Main Content ---
\section{Introduction}
\lipsum[1] % Generates one paragraph of lorem ipsum

\section{Methodology}
\subsection{Overview}
\lipsum[2] 

\section{Data Details}
\lipsum[3]

Table \ref{tab:example} highlights a generic categorisation of data points.

% --- Table Example using Booktabs ---
\begin{table}[h]
    \centering
    \caption{Generic data categories.}
    \label{tab:example}
    \begin{tabular}{lllr}
        \toprule
        \textbf{Category} & \textbf{Subcategory} & \textbf{Description} & \textbf{Count} \\
        \midrule
        Alpha & A1 & Primary item & 1539 \\
        \midrule
        Beta & B1 & Secondary item & 1510 \\
             & B2 & Tertiary item & 332 \\
             & B3 & Quaternary item & 2 \\
        \midrule
        Gamma & G1 & Quinary item & 1847 \\
              & G2 & Senary item & 269 \\
        \bottomrule
    \end{tabular}
\end{table}

\end{document}